\section{Restriction and unit-circle coefficients}
\label{sec:projective_restriction}

\begin{definition}[Pullback of scalar cohomology]
    \label{def:mpug_h2_pullback}
    \lean{ScalarCocycle.comap,
        ScalarCocycle.comap_apply,
        ScalarCocycle.mapCocycles₂_cocyclesOfIsMulCocycle₂,
        ScalarCocycle.map_H2π}
    \leanok
    \uses{def:h2_cx, def:scalar_representation_res}
    Let $f:K\to G$ be a group homomorphism. For a scalar $2$-cochain
    $\omega:G\times G\to\C^\times$, define
    $(f^*\omega)(a,b)=\omega(f(a),f(b))$.
    The restriction map
    $f^*:H^2(G,\C^\times)\to H^2(K,\C^\times)$ induced by $f$, with the
    identity of $\C^\times$ as coefficient map, satisfies
    $f^*[\omega]=[f^*\omega]$ for every $2$-cocycle $\omega$.
\end{definition}

\begin{definition}[Restriction of the trivial coefficients]
    \label{def:scalar_representation_res}
    \lean{ScalarCohomology.scalarRepresentationRes}
    \leanok
    \uses{lem:scalar_cohomology_trivial_action}
    Let $f:K\to G$ be a group homomorphism. Restricting the trivial
    representation of $G$ on $\C^\times$ along $f$ gives the trivial
    representation of $K$ on $\C^\times$, and the identity of $\C^\times$ is a
    morphism of representations between them. It is the coefficient map of
    the restriction of cochains and of cohomology along $f$.
\end{definition}

\begin{lemma}[Pullback preserves cocycles and cohomology]
    \label{thm:mpug_h2_pullback_preserves}
    \lean{ScalarCocycle.isMulCocycle₂_comap,
        ScalarCohomology.isMulCoboundary₂_div_comap}
    \leanok
    \uses{def:mpug_h2_pullback, def:cohomologous_to, def:is_cocycle}
    Let $f:K\to G$ be a group homomorphism. If $\omega$ is a scalar
    $2$-cocycle on $G$, then $f^*\omega$ is a scalar $2$-cocycle on $K$.
    If $\omega_1$ and $\omega_2$ are cohomologous scalar $2$-cochains on $G$,
    then $f^*\omega_1$ and $f^*\omega_2$ are cohomologous on $K$.
\end{lemma}

\begin{proof}\leanok
    \uses{def:mpug_h2_pullback, def:cohomologous_to, def:is_cocycle}
    For $a,b,c\in K$, substituting $f(a),f(b),f(c)$ into the cocycle equation
    for $\omega$ and using $f(ab)=f(a)f(b)$ and $f(bc)=f(b)f(c)$ shows that
    $f^*\omega$ is a $2$-cocycle.
    If $\omega_1(g,h)=\varphi(g)\varphi(h)\varphi(gh)^{-1}\omega_2(g,h)$,
    then with $\eta=\varphi\circ f$, substituting $g=f(a)$ and $h=f(b)$ gives
    $(f^*\omega_1)(a,b)=\eta(a)\eta(b)\eta(ab)^{-1}(f^*\omega_2)(a,b)$,
    so $\eta$ gives the required cohomology relation.
\end{proof}

\begin{lemma}[Exact extension of a subgroup cocycle]
    \label{thm:mpug_h2_exact_extension}
    \lean{ScalarCocycle.H2π_mem_range_map_subtype_iff,
        ScalarCocycle.isMulCocycle₂_coboundary_mul}
    \leanok
    \uses{def:mpug_h2_pullback, def:is_cocycle}
    Let $H\le G$ be any subgroup, let $\iota:H\hookrightarrow G$ be its
    inclusion, and let $\psi:H\times H\to\C^\times$ be a $2$-cocycle.
    Then $[\psi]$ belongs to the image of $\iota^*$ if and only if there is
    a $2$-cocycle $\Psi:G\times G\to\C^\times$ with
    $\Psi(a,b)=\psi(a,b)$ for every $a,b\in H$.
    No normalization of either cocycle is required.
    This is the adjustment of representatives implicit in the extension
    condition of \cite[lines~7042--7064]{FrancoRubio2025MPUGauging},
    here stated with $\C^\times$ coefficients.
\end{lemma}

\begin{proof}\leanok
    \uses{def:mpug_h2_pullback, def:cohomologous_to, def:is_cocycle}
    If $[\psi]=\iota^*[\omega]$, there is a cochain
    $\varphi:H\to\C^\times$ such that
    $\psi(a,b)=\varphi(a)\varphi(b)\varphi(ab)^{-1}\omega(a,b)$.
    Extending $\varphi$ to a function $\Phi:G\to\C^\times$ that takes value $1$
    outside $H$, define
    $\Psi(g,h)=\Phi(g)\Phi(h)\Phi(gh)^{-1}\omega(g,h)$.
    In the cocycle equation for $\Psi$, cancellation leaves the common
    factor $\Phi(g)\Phi(h)\Phi(k)\Phi(ghk)^{-1}$ multiplied by the two sides
    of the cocycle equation for $\omega$. Hence $\Psi$ is a cocycle, and
    $\Phi|_H=\varphi$ gives $\Psi|_{H\times H}=\psi$ exactly.
    Conversely, any exact extension $\Psi$ satisfies
    $\iota^*[\Psi]=[\psi]$.
\end{proof}

\begin{lemma}[Circle and nonzero-complex coefficients in degree two]
    \label{thm:mpug_circle_complex_units_h2}
    \lean{Complex.div_toUnits_unitsPhase,
        ScalarCocycle.circlePhase_isMulCocycle₂,
        ScalarCocycle.circlePhaseInclusion_isMulCocycle₂,
        ScalarCocycle.H2π_circlePhaseInclusion,
        ScalarCocycle.circlePhaseMathlibCocycle,
        ScalarCocycle.H2π_map_circlePhaseMathlibCocycle,
        ScalarCohomology.circleToComplexUnits_map_injective,
        ScalarCohomology.circleToComplexUnits_map_surjective,
        ScalarCohomology.circleToComplexUnits_map_bijective}
    \leanok
    \uses{thm:mpug_h2_group_cohomology_comparison,
        thm:mpug_positive_generalized_cocycle}
    Let $G$ be finite, with $U(1)$ and $\C^\times$ carrying the trivial
    $G$-actions.  Inclusion of the unit circle induces an equivalence
    \begin{align}
        H^2(G,U(1))
         & \simeq H^2(G,\C^\times).
        \label{eq:mpug_circle_complex_units_h2}
    \end{align}
    More concretely, for a $\C^\times$-valued cocycle $\omega$, define its
    pointwise phase and radial residual by
    \begin{align}
        u(g,h)
         & =\frac{\omega(g,h)}{|\omega(g,h)|}\in U(1),
         &
        r(g,h)
         & =\frac{\omega(g,h)}{u(g,h)}=|\omega(g,h)|.
        \label{eq:mpug_circle_phase_residual}
    \end{align}
    Then $u$ is a cocycle, $r$ is a positive coboundary, and hence $\omega$
    is cohomologous to the inclusion of $u$.  The inclusion map is injective
    for every group because pointwise phase is its left inverse; finiteness is
    used only to prove surjectivity.

    As recorded in \cite{gap:fbc25_circle_complex_units_cohomology}, this
    algebraic group-cohomology statement for a finite discrete group replaces
    the comparison of sheaf cohomology with topological coefficient groups
    invoked in the footnote at
    \cite[line~1931]{FrancoRubio2025MPUGauging}.
\end{lemma}

\begin{proof}\leanok
    \uses{thm:mpug_positive_generalized_cocycle}
    The phase map $q(z)=z/|z|$ is multiplicative and satisfies
    $q(u)=u$ for $u\in U(1)$, so it preserves the cocycle equation and
    retracts inclusion.  The polar identity in
    (\ref{eq:mpug_circle_phase_residual}) identifies the residual $r$ with a
    positive cocycle.  Applying
    Lemma~\ref{thm:mpug_positive_generalized_cocycle} to the one-point label
    set writes $r=d\chi$.  Therefore $\omega\sim u$, which proves
    surjectivity of coefficient inclusion on degree-two classes.  Its left
    inverse induced by $q$ proves injectivity.
\end{proof}
